% ============================================================================
% Build: tectonic main.tex            Tables: python3 tables.py  (writes tables.tex)
% Every number in this paper is generated from cached measurements by tables.py
% or quoted from a named file; see ../figures/README.md for the commands.
%
% FRAMING NOTE (2026-09-29): no novelty claim. Emitting one neutral feature
% representation into several CAD systems was done by the macro-parametric line
% (Mun/Han 2003; Proficiency) and, contemporaneously, by CADIR (arXiv 2608.00891,
% 2026-08-01). See ../prior-art/SURVEY-multivendor.md. The contribution is an
% open specification, measurements, and what cross-backend disagreement found.
% ============================================================================
\documentclass[11pt]{article}

\usepackage[margin=1.1in]{geometry}
\usepackage[T1]{fontenc}
\usepackage[utf8]{inputenc}
\usepackage{lmodern}
\usepackage{microtype}
\usepackage{amsmath}
\usepackage{amssymb}
\usepackage{booktabs}
\usepackage{graphicx}
\usepackage{siunitx}
\usepackage{tikz}
\usetikzlibrary{positioning,arrows.meta}
\usepackage[numbers,sort&compress]{natbib}
\usepackage[hidelinks]{hyperref}
\usepackage{xspace}
\usepackage{listings}
\lstset{basicstyle=\ttfamily\footnotesize,columns=fullflexible,frame=single,breaklines=true}

\newcommand{\code}[1]{\texttt{\small #1}}
\newcommand{\VERIFIED}{\textsc{verified}}
\newcommand{\FCverified}{10}
\newcommand{\FCnote}{The exception is FTC-06 (build fails); \S\ref{sec:defects} traces the cause.}
\newcommand{\NatOnshape}{11}
\newcommand{\NatCadmpeg}{1}
   % measured numbers, generated by tables.py

\title{\textbf{The \code{featuretree} IR: an Open Feature-Tree Format,\\
Emitted Natively into Several CAD Systems and Checked by Re-Reading}}
\author{Daniel Newcome\\[2pt]
  \small punkfab\\
  \small\url{https://github.com/punkfab/featuretree}}
\date{October 2026 (version 2)}

\begin{document}
\maketitle

\begin{abstract}
When a program or a language-model agent writes a part, a person still has to be
able to open it, see how it was built, and change one dimension without starting
again. That needs a shared, small description of the feature tree --- the
``timeline'' every history-based CAD system shows --- that an agent can emit, each
CAD system can rebuild natively, and an edit made in any of them can be read
back into. Feature-based modelling has been studied since the 1980s and no
complete vocabulary exists; we publish a deliberately small one and a rule for
growing it. The \code{featuretree} IR is a declarative JSON list of named
features --- sketches with arcs, pad, pocket, fillet, revolve and placed cuts,
with draft --- whose references to faces and edges are stored as queries resolved
against live geometry, never as kernel identifiers. One IR is emitted as a
native, editable tree in FreeCAD, build123d and Onshape, and as scripts for
Autodesk Fusion and SolidWorks. Every emitted part is checked the same way: STEP
export, a volume and Boolean-free IoU gate against a reference build, and a
per-feature volume trace that names the first feature where two backends
disagree. We report the measurements, including the failures. On the eleven
trees recovered from the NIST MBE PMI corpus FreeCAD reproduces the reference on
\FCverified{} of 11; Onshape, limited by its free plan's API quota, has run
\OSrun{} and passes \OSverified{}. An edit made in the FreeCAD tree, read back by
feature name and rebuilt by the reference, agrees on \NudgeAgree{} of
\NudgeTried{} single-dimension edits; before a defect that this measurement found
was fixed, it agreed on \NudgeBeforeAgree{}. Of the \VocabTotal{} features in the
designers' own trees of the same parts, \VocabKindPct\% have a kind of their own in
the IR; three more kinds (hole, pattern, chamfer) would bring that to
\VocabNextPct\%. Running the backends against each other found six defects,
including one in the reference backend and two rules the specification had never
stated. Neutral feature exchange is not new --- it goes back to the
macro-parametric work of 2003, and a closely related system appeared this year
--- so we claim no novelty. We offer an open specification, an MIT-licensed
implementation, and measurements that can be reproduced.
\end{abstract}

% ----------------------------------------------------------------------------
\section{Introduction}

Moving a parametric part between CAD systems loses its design history. STEP,
the neutral exchange standard, has carried shape reliably for decades; its parts
for construction history (ISO~10303-55/108/111) were specified but, as their own
authors reported, not adopted by commercial translators~\cite{TODO-pratt-kim}.
The alternatives that did ship --- neutral modelling commands replayed through
each system's API~\cite{TODO-mun-han,safdar2020macroparametric}, and commercial
bridges such as Proficiency~\cite{TODO-rappoport} --- showed that the approach
works, and identified the hard problems: feature mapping, persistent naming and
constraint translation~\cite{safdar2020macroparametric}.

The question has come back for a new reason. Parts are now written by programs
and by language-model agents, and most of that output is either code for one
kernel (CadQuery, FreeCAD Python) or bare geometry~\cite{alrashedy2025cadcodeverify,rukhovich2025cadrecode,wu2021deepcad}.
A person reviewing such a part needs to nudge it: open it in the system they
know, find the feature, change a dimension, and hand the result back. Bare
geometry does not allow that, and code for one kernel allows it only to someone
who reads that code. What is needed is a representation of the feature tree
itself that is small enough for a program to emit, explicit enough for several
CAD systems to rebuild natively, checkable, and able to take an edit back from
any of them. Which features such a representation should contain is an old
question: feature-based modelling was surveyed thirty years
ago~\cite{shah1991assessment,salomons1993review}, the vendors' vocabularies
still differ, and no list is complete. We do not try to settle it. We fix a
small subset, measure how far it reaches, and state how it grows
(\S\ref{sec:grow}). CADIR, which
appeared in August 2026, targets the same need with an executable Python IR and
adapters for FreeCAD, SolidWorks and Fusion~\cite{liu2026cadir}.
\code{featuretree} was developed independently (public since June 2026) and makes
different choices; we describe them and measure the result.

\paragraph{Contributions.}
\begin{enumerate}
\item \textbf{A specification} (\S\ref{sec:ir}): a declarative, non-executable
  JSON feature list with seven feature kinds, query-based references, fixed
  coordinate and sign conventions, and well-formedness rules with a validator.
\item \textbf{Five backends} (\S\ref{sec:backends}): build123d (the reference),
  FreeCAD and Onshape, which we run and measure, and generated scripts for
  Fusion and SolidWorks, which we have not yet run.
\item \textbf{One verification method for every backend}
  (\S\ref{sec:verify}): STEP export, a volume and Boolean-free IoU gate, a
  per-feature volume trace, and read-back of native files through two
  independent decoders.
\item \textbf{A mapping to each system's timeline} (\S\ref{sec:timeline}): what
  every IR kind is called in each tree, which dimensions read back by name, and
  a measured edit round trip.
\item \textbf{Measurements} (\S\ref{sec:results}): the backends on the eleven
  recovered NIST trees and on coverage parts; the six defects that cross-backend
  disagreement exposed; the edit round trip; a two-decoder study of NIST's
  native SolidWorks files; and how much of the designers' own feature vocabulary
  the IR covers.
\item \textbf{A rule for growing the vocabulary} (\S\ref{sec:grow}).
\end{enumerate}

% ----------------------------------------------------------------------------
\section{The IR}\label{sec:ir}

A part is \code{\{"name": ..., "features": [...]\}}: an ordered list of named
features. Order is construction order; a feature may refer only to earlier
features, and only by name. Table~\ref{tab:kinds} lists the kinds.
Listing~\ref{lst:ir} is a complete part.

\begin{table}[t]
\centering\small
\begin{tabular}{@{}lp{0.62\linewidth}@{}}
\toprule
Kind & Parameters and meaning \\
\midrule
\code{sketch} & \code{circles} $(c_x,c_y,r)$, \code{rects} $(w,h,c_x,c_y)$, \code{polys}
  (closed loops of $(x,y[,b])$, $b$ = bulge); on plane \code{XY} or \code{XZ}, or
  \code{on} a face given by the query \code{\{"face\_of", "side": top|bottom\}} \\
\code{pad} & extrude a sketch: \code{length}, \code{symmetric}, \code{taper} (deg) \\
\code{pocket} & cut a sketch: \code{through} or blind \code{length}, \code{taper} \\
\code{fillet} & \code{radius} on the edges selected by a query, e.g.\ the largest circle on the top \\
\code{revolve} & revolve an \code{XZ} profile about $Z$ by \code{angle} \\
\code{polar\_pocket} & \code{count} cylinders tangent to a circle (roller pockets) \\
\code{prism\_cut} & cut a profile drawn on the plane (\code{origin}, \code{normal}, \code{xdir}), \code{depth} along the normal, \code{taper} \\
\bottomrule
\end{tabular}
\caption{Feature kinds. Lengths in millimetres, angles in degrees.}
\label{tab:kinds}
\end{table}

\begin{lstlisting}[caption={A complete part: a plate, a pocket in its top face, and a fillet.},label={lst:ir},float=t]
{"name": "plate", "features": [
  {"kind": "sketch", "name": "outline", "plane": "XY",
   "polys": [[[-20,-15],[20,-15,0.4142],[20,15],[-20,15]]]},
  {"kind": "pad", "name": "body", "sketch": "outline", "length": 10},
  {"kind": "sketch", "name": "recess_sk", "on": {"face_of": "body", "side": "top"},
   "circles": [[0, 0, 6]]},
  {"kind": "pocket", "name": "recess", "sketch": "recess_sk", "through": false, "length": 4},
  {"kind": "fillet", "name": "round", "radius": 1, "select": {"circles": "top_outer"}}]}
\end{lstlisting}

\paragraph{References are queries.} A face-attached sketch names ``the top face''
and a fillet names ``the largest circular edge on top'', and every backend
resolves the query against its own live geometry each time the part is built.
This is the approach of generic naming~\cite{capoyleas1996generic} and of
geometry-based semantic identifiers~\cite{wang2005geometry}, and of the selectors
in CadQuery, build123d and Onshape's FeatureScript; it avoids tracking kernel
names through history~\cite{kripac1997mechanism,marcheix2002survey}. The known
cost applies: after a large edit a query can resolve to a different entity.

\paragraph{Conventions.} Coordinates of a face-attached sketch stay global:
a circle at $(0,0)$ on the top face is above the origin, whatever frame the
backend gives that face. \code{XZ} maps $(u,v)$ to $(u,0,v)$. A pad grows
$+Z$, or $-Z$ from a bottom face. A bulge is the DXF arc factor $\tan(\theta/4)$,
its arc's midpoint on the chord's left normal, and $|b|<10^{-4}$ is a straight
line. A positive taper shrinks the profile along the extrusion (for a pocket,
as it goes deeper). A sketch's region is its loops under the nesting rule: a
loop inside another is a hole, an island in a hole is material again, and
separate loops are separate regions.

\paragraph{Well-formedness.} \code{ir.validate(spec)} checks what every backend
is entitled to assume: unique names; references to earlier sketches; closed
loops (at least three segments, or two if one is an arc) with no zero-length
segment and no self-intersection; within one sketch, loop boundaries that never
touch or cross; a taper only on a blind pocket; an orthogonal placement frame.
Two of these rules (non-touching loops and the bulge threshold) were added after
backends disagreed about parts that broke them (\S\ref{sec:defects}).

% ----------------------------------------------------------------------------
\section{Backends}\label{sec:backends}

\begin{table}[t]
\centering\small
\begin{tabular}{@{}llll@{}}
\toprule
Backend & Driven by & References become & Status \\
\midrule
build123d & in-process Python (OCCT) & queries on the live solid & reference \\
FreeCAD 1.0 & headless \code{freecadcmd}, PartDesign & \code{FaceN}/\code{EdgeN} re-resolved per build & measured \\
Onshape & REST v6 feature JSON & FeatureScript query strings & measured \\
Fusion & generated script (\code{adsk.fusion}) & queries evaluated in the script & not yet run \\
SolidWorks & generated script (COM, pywin32) & queries evaluated in the script & not yet run \\
\bottomrule
\end{tabular}
\caption{Backends. ``Measured'' means built and gated in this paper.}
\label{tab:backends}
\end{table}

Table~\ref{tab:backends} summarises them. \textbf{build123d} is the reference:
every other backend is compared with its build of the same IR. \textbf{FreeCAD}
writes a native PartDesign body in which each feature's label is its IR name, so
a dimension edited in the GUI reads back by name. \textbf{Onshape} receives
native Sketch, Extrude, Fillet, Revolve and Plane features. Every reference is
a FeatureScript query string rather than an entity identifier: a profile is
\code{qContainsPoint(qSketchRegion(id), p)} for an interior point $p$ of each
material region, a face-attached sketch plane is the planar $Z$-normal face
farthest along $\pm Z$, and a fillet's edges are the circular edges containing
points computed from the same rule the reference uses. Onshape re-evaluates them
on every regeneration, so the tree stays editable. Placed cuts on oblique planes
and polar pockets are not yet emitted. \textbf{Fusion} and \textbf{SolidWorks}
cannot be driven from outside the program, so the emitter writes a
self-contained script: the IR as data, the reference volume after every
feature, and a small runtime that re-authors each feature and resolves each
query inside the program, then exports STEP and a report. The shared geometry
those runtimes draw from is tested by replaying it with build123d, but the
scripts have not yet run in either program, so we report no results for them.

We looked for a way to write a SolidWorks file without SolidWorks. The
open-source decoder we use to read \code{.sldprt} files, cadmpeg, also lists
\code{.sldprt} and \code{.f3d} as write targets. At version 0.6.0 its writers
decline a plain box read from STEP (``SLDPRT semantic writer does not encode
explicit edge parameter ranges''; for Fusion, a file with no native source
cannot be serialised), the same refusal is in its source on 4 October 2026, and
none of its 24 public forks changes the writer. So today the open route reads
SolidWorks files but does not write them, and a SolidWorks result still needs
the program.

% ----------------------------------------------------------------------------
\section{The IR as a shared timeline}\label{sec:timeline}

Each history-based CAD system shows the user an ordered list of features: the
model tree in FreeCAD, the feature list in Onshape, the timeline in Fusion, the
FeatureManager tree in SolidWorks. The IR is that list with the vendor removed.
Table~\ref{tab:timeline} gives the correspondence the backends implement. Every
emitted feature carries its IR name as its name in the target's tree, and that
name is the only key used in either direction.

\begin{table}[t]
\centering\footnotesize
\begin{tabular}{@{}lp{0.17\linewidth}p{0.2\linewidth}p{0.2\linewidth}l@{}}
\toprule
IR kind & FreeCAD & Onshape & Fusion, SolidWorks & Reads back \\
\midrule
\code{sketch} & Sketch & Sketch & Sketch & circle radii \\
\code{pad} & Pad & Extrude (new/add) & Extrude (join), Boss-Extrude & length \\
\code{pocket} & Pocket & Extrude (remove) & Extrude (cut), Cut-Extrude & length \\
\code{fillet} & Fillet & Fillet & Fillet & radius \\
\code{revolve} & Revolution & Revolve & Revolve & angle \\
\code{prism\_cut} & plane, sketch, Pocket & plane, sketch, Extrude (remove) & plane, sketch, cut & depth \\
\code{polar\_pocket} & $n$ subtractive cylinders & not emitted & $n$ placed cuts & -- \\
\bottomrule
\end{tabular}
\caption{What each IR kind becomes in each system's tree, and the driving
dimension that is read back by feature name. The Fusion and SolidWorks columns
describe the generated scripts, which have not run.}
\label{tab:timeline}
\end{table}

\paragraph{The edit loop.} A person changes a dimension in the target's own
interface. A reader for that target returns \code{\{feature name: \{parameter:
value\}\}}; \code{ir.update\_from\_params} writes those values into the IR; every
other backend rebuilds from the IR. Nothing else flows back: a feature added,
deleted or reordered in the target is not read, and sketch geometry other than
circle radii is not read. Readers exist for FreeCAD and Onshape, and inside the
Fusion and SolidWorks scripts. \S\ref{sec:nudge} measures the FreeCAD loop.

\paragraph{What the table hides.} The mapping is one-to-one only for the first
five rows. A placed cut becomes three tree entries in every target, and a polar
pocket becomes $n$; a person reading the tree sees the parts, not the IR
feature. The reverse problem is larger: a designer's hole or pattern feature has
no kind of its own and arrives as plain cuts (\S\ref{sec:vocab}).

% ----------------------------------------------------------------------------
\section{Verification}\label{sec:verify}

A backend's output is exported to STEP and compared with the reference build of
the same IR. It passes (\VERIFIED{}) when volumes agree within
\SI{0.5}{\percent} and the intersection-over-union is at least
\SI{99.5}{\percent}. IoU is estimated by classifying uniformly sampled points
against both solids~\cite{tilove1980smc}, never by a Boolean, because
OpenCASCADE's Booleans returned wrong unions on these parts; robustness failures
of this kind are well documented~\cite{kettner2008robustness,zhou2016mesharr}.
The volume check is the idea behind CAx-IF's geometric validation
properties~\cite{caxif_gvp_2014} and Sap and Shapiro's invariant-property
testing~\cite{sap2019invariant}; IoU additionally catches misplaced material
that leaves volume unchanged.

Each backend also records the volume after every feature. When a part fails,
the first feature whose volume leaves the reference is where to look; every
defect in \S\ref{sec:defects} was found that way.

For native files (\code{.sldprt}, \code{.f3d}) a third check applies:
the file is decoded to STEP by two independent decoders, and its stored feature
history is matched to the IR by name, kind and driving value. Comparing two
implementations of the same function is differential
testing~\cite{mckeeman1998differential,yang2011csmith}; it answers the oracle
problem~\cite{barr2015oracle} that a single decoder's own consistency check
cannot (\S\ref{sec:native}).

% ----------------------------------------------------------------------------
\section{Results}\label{sec:results}

\subsection{FreeCAD and Onshape against the reference}

The corpus is the eleven feature trees that \code{featuretree}'s recogniser
recovers from the NIST MBE PMI test parts~\cite{lipman2017nistpmi}, 2 to 144
features each (the recogniser and its accuracy are the subject of a companion
paper; here the trees are simply inputs). Table~\ref{tab:corpus} gives
FreeCAD's result on every part and Onshape's where it ran.

\begin{table}[t]
\centering\small
\begin{tabular}{@{}lrrlrrlrr@{}}
\toprule
& & & \multicolumn{3}{c}{FreeCAD} & \multicolumn{3}{c}{Onshape} \\
\cmidrule(lr){4-6}\cmidrule(l){7-9}
Part & Features & Valid & Verdict & IoU \% & $|\Delta V|$ \% & Verdict & IoU \% & $|\Delta V|$ \% \\
\midrule
CTC-01 & 54 & yes & \VERIFIED{} & 100.00 & 0.00 & 53/54 built & 99.36 & 0.58 \\
CTC-02 & 144 & yes & \VERIFIED{} & 99.90 & 0.05 & not run &  &  \\
CTC-03 & 64 & no & \VERIFIED{} & 100.00 & 0.00 & not run &  &  \\
CTC-04 & 108 & yes & \VERIFIED{} & 100.00 & 0.00 & not run &  &  \\
CTC-05 & 57 & yes & \VERIFIED{} & 99.78 & 0.23 & not run &  &  \\
FTC-06 & 27 & no & build fails &  &  & not run &  &  \\
FTC-07 & 29 & yes & \VERIFIED{} & 100.00 & 0.00 & not run &  &  \\
FTC-08 & 37 & yes & \VERIFIED{} & 100.00 & 0.00 & not run &  &  \\
FTC-09 & 80 & yes & \VERIFIED{} & 100.00 & 0.00 & not run &  &  \\
FTC-10 & 38 & yes & \VERIFIED{} & 100.00 & 0.00 & not run &  &  \\
FTC-11 & 2 & yes & \VERIFIED{} & 100.00 & 0.00 & \VERIFIED{} & 100.00 & 0.00 \\
\bottomrule
\end{tabular}
\caption{Each recovered NIST feature tree built by FreeCAD and Onshape and gated against the build123d reference build of the same IR (\VERIFIED{}: IoU $\geq$ 99.5\% and $|\Delta V| \leq$ 0.5\%). ``Valid'': passes \code{ir.validate}. Onshape's run stopped at the free plan's daily API quota.}
\label{tab:corpus}
\end{table}

\begin{table}[t]
\centering\small
\begin{tabular}{@{}llrrr@{}}
\toprule
Part & Feature & Onshape (mm$^3$) & Reference (mm$^3$) & Difference \\
\midrule
plate (sample) & \code{body} & 11,497.345 & 11,497.345 & +0.000 \\
 & \code{hole} & 11,497.345 & 11,497.345 & +0.000 \\
every kind & \code{base} & 22,028.273 & 22,028.273 & +0.000 \\
 & \code{top\_pk} & 21,784.056 & 21,788.231 & -4.175 \\
 & \code{hub} & 24,328.746 & 24,332.921 & -4.175 \\
 & \code{side\_slot} & 24,157.241 & 24,161.417 & -4.176 \\
 & \code{thru} & 24,006.445 & 24,010.621 & -4.176 \\
 & \code{round} & 23,980.157 & 23,984.332 & -4.175 \\
taper A & \code{base} & 12,000.000 & 12,000.000 & +0.000 \\
 & \code{up} & 12,432.802 & 12,432.802 & +0.000 \\
 & \code{down} & 12,668.025 & 12,668.025 & +0.000 \\
 & \code{sym} & 12,724.573 & 12,724.573 & +0.000 \\
taper B & \code{base} & 12,000.000 & 12,000.000 & +0.000 \\
 & \code{pk\_bottom} & 11,826.944 & 11,826.944 & +0.000 \\
 & \code{pk\_top} & 11,632.467 & 11,632.467 & +0.000 \\
 & \code{side} & 11,550.769 & 11,550.769 & +0.000 \\
\bottomrule
\end{tabular}
\caption{Onshape coverage parts: volume after every solid feature against the build123d reference. The $-4.175$ offset enters at \code{top\_pk} (a pocket crossing a concave arc edge) and is carried through; everything else agrees to three decimals.}
\label{tab:coverage}
\end{table}

\begin{table}[t]
\centering\small
\begin{tabular}{@{}lrrrlrr@{}}
\toprule
& \multicolumn{2}{c}{Onshape import} & \multicolumn{2}{c}{cadmpeg 0.6.0} & \multicolumn{2}{c}{Solid operations} \\
\cmidrule(lr){2-3}\cmidrule(lr){4-5}\cmidrule(l){6-7}
Part & IoU \% & $\Delta V$ \% & IoU \% & STEP & Designer & Recovered \\
\midrule
CTC-01 & 100.00 & -0.000 & 21.39 & invalid & 10 & 39 \\
CTC-02 & 99.94 & +0.009 & 3.47 & invalid & 42 & 89 \\
CTC-03 & 100.00 & -0.164 & 11.14 & invalid & 15 & 33 \\
CTC-04 & 100.00 & -0.001 & 22.37 & no solid & 24 & 93 \\
CTC-05 & 100.00 & -0.000 & 6.99 & no solid & 17 & 40 \\
FTC-06 & 100.00 & +0.000 & 25.83 & invalid & 21 & 18 \\
FTC-07 & 99.93 & +0.009 & 0.00 & no solid & 28 & 21 \\
FTC-08 & 100.00 & +0.008 & 6.95 & invalid & 30 & 32 \\
FTC-09 & 100.00 & -0.000 & 80.72 & invalid & 33 & 40 \\
FTC-10 & 100.00 & +0.001 & 26.46 & no solid & 19 & 25 \\
FTC-11 & 100.00 & -0.000 & 100.00 & valid & 1 & 1 \\
\bottomrule
\end{tabular}
\caption{NIST's native SolidWorks 2018 files decoded to STEP by two decoders and compared with NIST's own STEP of each part; and the number of solid operations in the designer's tree (as decoded) against the tree \code{featuretree} recovers from the STEP.}
\label{tab:native}
\end{table}


FreeCAD reproduces the reference on \FCverified{} of 11 parts. \FCnote{}

Onshape has run on \OSrun{} parts and is \VERIFIED{} on \OSverified{}, FTC-11.
On CTC-01 it builds 53 of 54 features; one placed cut, whose profile edge lies
exactly on the part's outer wall, is rejected by Onshape's kernel, and the part
reaches \SI{99.36}{\percent} IoU, just under the gate. On CTC-02, the largest
tree, it builds 93 of 144: ten of the stacked pads that make up the body
regenerate with an error, most of the later cuts then have nothing to cut, and
the result is far from the reference. We have not yet found why those pads fail.
The free plan's daily API quota (about 1,500 calls) is what limits this column:
our emitter spends several calls on each feature, CTC-02 alone used a day's
quota, and the remaining eight parts have not run. Cutting the calls per feature
is the next change to that backend.

Onshape was also run on coverage parts built to exercise every path it supports
(Table~\ref{tab:coverage}). With the conventions pinned, every feature's volume
matches the reference to three decimals, with one exception: a pocket whose
circle crosses a concave arc edge removes \SI{4.175}{mm^3} more in Onshape than
the exact answer, computed independently, which build123d matches. We have not
explained it.

\subsection{What disagreement found}\label{sec:defects}

Running backends against each other is a test of each of them and of the
specification. The trace located six defects, none of which any single backend
reported:

\begin{enumerate}
\item \textbf{The reference was wrong.} build123d decided that a sketch was on
  the bottom face by its height being below zero, so a pad from a bottom face at
  $z=0$, the common case, grew upward into the part. Onshape disagreed on the
  first such feature. The rule now uses the sketch's declared side. No published
  result changed: all eleven recovered trees give identical reference volumes
  before and after.
\item \textbf{FreeCAD kept one solid.} A PartDesign body is one solid by
  default, so a sketch of several separate outlines padded only one of them,
  exactly halving the first feature of CTC-03 and CTC-04. FreeCAD 1.0's
  \code{AllowCompound} restores the specified union; CTC-03 went from
  \SI{2.5}{\percent} IoU to \VERIFIED{}.
\item \textbf{A rule the specification lacked: loops must not touch.} Two
  recovered trees contain sketches whose loops share boundary. build123d unions
  loops one at a time and accepts them; FreeCAD rejects some (FTC-06).
  \code{ir.validate} now reports them; the recogniser should not produce them.
\item \textbf{A rule the specification lacked: near-zero bulges.} CTC-02 contains
  arcs of bulge $5\times10^{-6}$, radius about two kilometres. FreeCAD's sketch
  solver fails on them. Below $10^{-4}$ a bulge now means a straight line in
  every backend; the reference volume of CTC-02 moves by
  \SI{-0.000009}{\percent}.
\item \textbf{Onshape's draft direction} is the opposite of our first guess;
  a tapered boss grew where it should shrink. Pinned by matching the reference on
  tapered pads and pockets in both directions.
\item \textbf{Failure behaviour differs.} A cut that removes nothing is an error
  in Onshape and silently does nothing in build123d; the trace shows both.
\end{enumerate}

\subsection{Reading native files back: one decoder is not enough}\label{sec:native}

NIST publishes the same test parts as native SolidWorks 2018 files. We decoded
each to STEP with an open-source clean-room decoder, cadmpeg 0.6.0, and with
Onshape's commercial importer, and compared both with NIST's own STEP of the
part (Table~\ref{tab:native}). Onshape matches on \NatOnshape{} of 11 (IoU
$\geq$ \SI{99.93}{\percent}); cadmpeg on \NatCadmpeg{}. Four of cadmpeg's
wrong parts pass its own model check and convert without complaint; four more
yield a shell that is not a solid. Its decoding of the feature history, on the
other hand, looked right on the same files. The practical rule we adopted: a
native file's geometry is judged by the commercial importer when available, its
history by the open decoder. We reported the geometry failures upstream with a
reproduction.

\subsection{An edit, round the loop}\label{sec:nudge}

For each part FreeCAD could build quickly (the two samples and eight NIST trees;
CTC-02 and CTC-04 take minutes per rebuild and were left out) we took up to six
editable dimensions spread through the tree. Each was set to 1.1 times its value
in the FreeCAD file, alone, as a person would in the GUI; FreeCAD recomputed;
every parameter was read back into the IR; and the reference rebuilt from the
IR (Table~\ref{tab:nudge}).

\NudgeAgree{} of \NudgeTried{} edits agree. The first run agreed on only
\NudgeBeforeAgree{} of \NudgeBeforeTried{}: \NudgeBeforeLost{} edits to placed
cuts never reached the IR, because FreeCAD holds a placed cut as a Pocket and
reports its depth as \code{length}, while the IR reader looked for \code{depth}.
The edit was silently dropped and the rebuilt part was the old part. The same
run found that the function for editing a sketch circle called a FreeCAD method
that does not exist in version 1.0. Both are fixed, and both had been in the
code since the round trip was first written; no test had exercised them on a
recovered tree.

The three that still fail are all edits to the length of a base pad. On FTC-07
(two edits) FreeCAD's recomputed body is invalid, with a negative volume, while
the reference rebuilds a valid solid. On CTC-01 neither system produces a solid.
The recovered trees place later cuts at absolute coordinates rather than
relative to the pad they cut, which is the likely reason, but we have not
diagnosed these and report them as failures.

\subsection{How much of a designer's vocabulary}\label{sec:vocab}

The native SolidWorks files carry the designers' own feature trees, which cadmpeg
decodes. Table~\ref{tab:vocab} counts their solid-modifying features over the
eleven parts and says what the IR has for each. \VocabKind{} of \VocabTotal{}
(\VocabKindPct\%) are extrudes, fillets and revolves, which have kinds of their
own. Almost all of the rest are holes, patterns and chamfers
(\VocabNext{} features): the IR can hold their shape, as cuts, but not the fact
that the designer made one hole feature with a counterbore or one pattern of
twelve. Those three kinds would bring the share with a kind of its own to
\VocabNextPct\%. Two cautions apply. The count says a kind exists, not that every
option of the designer's feature is covered: the IR's fillet, for example, can
select only circular edges on the top face. And eleven parts from one test suite
modelled in one CAD system are a small sample of how people model.

The same gap shows in tree length. The trees our recogniser recovers from the
STEP files are 2 to 4 times longer than the designers' on the CTC parts (CTC-01:
39 operations against 10), mostly because one hole or pattern feature stands for
many of our cuts (Table~\ref{tab:native}).

% ----------------------------------------------------------------------------
\section{Growing the vocabulary}\label{sec:grow}

The subset is meant to grow, and \S\ref{sec:defects} shows what an
under-specified kind costs: two backends that each look correct and disagree. A
kind is added when all of the following hold.

\begin{enumerate}
\item \textbf{Its meaning is written without reference to any kernel}: the
  parameters, their units, every sign and direction convention, and what happens
  at the edges (a cut that removes nothing, a zero length).
\item \textbf{The validator checks what backends may assume} about it.
\item \textbf{The reference backend builds it}, and a second backend on a
  different kernel agrees with the reference at that feature, in the volume
  trace, on parts written to exercise each convention in both directions.
\item \textbf{It has a row in Table~\ref{tab:timeline}}: a native feature of the
  same meaning in each target, or an explicit statement that the target gets a
  decomposition.
\item \textbf{Its driving dimensions read back by name}, and the edit round trip
  of \S\ref{sec:nudge} passes for it.
\end{enumerate}

By this rule the next kinds are the ones Table~\ref{tab:vocab} ranks: hole (with
its counterbore, countersink and thread callouts as parameters), linear and
circular pattern of an earlier feature, and chamfer. Each has a native feature of
the same name in all four targets. A draft applied to existing faces, shell and
rib are rarer in this sample and wait. Sketch constraints are a separate
question that this vocabulary does not address.

% ----------------------------------------------------------------------------
\section{Related work}\label{sec:related}

\textbf{Neutral feature exchange.} STEP's construction-history parts
(ISO 10303-55/108/111/112)~\cite{TODO-pratt-kim,kim2008designintent,kim2011procedural2d};
the macro-parametric approach and its test rallies between commercial
systems~\cite{TODO-mun-han,safdar2020macroparametric,kim2019assembly};
Proficiency's universal product representation~\cite{TODO-rappoport}; BYU's
multi-user neutral database~\cite{TODO-shumway}. Interoperability's cost motivated
all of it~\cite{brunnermeier2002interop}. \textbf{Closest system:}
CADIR~\cite{liu2026cadir} records a construction graph from an executable Python
IR and rebuilds native histories in FreeCAD, SolidWorks and Fusion, matching
topology by geometric signature; \code{featuretree} differs in being a
declarative data format, in storing queries rather than matched selections, in
targeting Onshape and build123d, and in sharing its IR with a STEP recogniser.
\textbf{CAD as a language.} Code-CAD (OpenSCAD, CadQuery, build123d,
FeatureScript, KCL)~\cite{machado2019openscad} and programming-language
treatments~\cite{nandi2018functional,nandi2020szalinski,TODO-cascaval,ritchie2023neurosymbolic};
learned construction sequences, mostly sketch and extrude~\cite{seff2020sketchgraphs,wu2021deepcad,xu2022skexgen,TODO-willis-fusion,khan2024text2cad,rukhovich2025cadrecode}.
\textbf{Persistent naming}~\cite{capoyleas1996generic,kripac1997mechanism,bidarra2002persistent,marcheix2002survey,wang2005geometry,farjana2018mechanisms}.
\textbf{Language models for CAD} emit CadQuery, FreeCAD Python or FeatureScript
for one system each~\cite{makatura2023llmdesign,alrashedy2025cadcodeverify,pyatov2026cadfs,hepworth2026iterative,mallis2025cadassistant}.
\textbf{Exchange quality}: validation properties~\cite{caxif_gvp_2014},
invariant-property and query-based interoperability~\cite{hoffmann2014queries,sap2019invariant,sap2025interop},
benchmarks~\cite{gerbino2004interop}, and NIST's conformance
testing~\cite{lipman2017nistpmi,lipman2020gdt}.

% ----------------------------------------------------------------------------
\section{Limitations}

\begin{itemize}
\item \textbf{Fusion and SolidWorks have not run.} Their scripts exist and their
  geometry is replay-tested, but API behaviour, sign conventions and profile
  selection are unverified until a run passes the gate.
\item \textbf{Onshape's corpus run is incomplete} (\OSrun{} of 11 parts) because of
  the free plan's API quota; one of the three fails badly for a reason we have
  not found; and the backend does not yet emit oblique placed cuts or polar
  pockets.
\item \textbf{The edit loop carries dimensions only.} Features added, removed or
  reordered in a target are not read back; the round trip was measured for
  FreeCAD only, on up to six edits per part, one at a time, each a 10\% change.
\item \textbf{The vocabulary is small.} About half of the features designers
  used on these parts arrive without their identity (\S\ref{sec:vocab}).
\item \textbf{The reference is OpenCASCADE, and so is FreeCAD.} Their agreement
  is partly self-consistency. Onshape (Parasolid) is the independent kernel.
\item \textbf{Single bodies, no sketch constraints, no assemblies} in this IR;
  sketches arrive in the target unconstrained.
\item \textbf{Queries can re-resolve wrongly} after large edits; we have not
  measured how often.
\item \textbf{The corpus is eleven parts} and the trees are machine-recovered,
  not hand-authored.
\end{itemize}

% ----------------------------------------------------------------------------
\section{Availability}

Code, specification (\code{ir.py}), validator and every backend are MIT-licensed
at \url{https://github.com/punkfab/featuretree}; the numbers here come from the commit tagged \code{ir-paper-v2}. The measurements are regenerated by:
\begin{itemize}
\item \code{paper/figures/freecad\_corpus.py} and \code{onshape\_corpus.py} (Table~\ref{tab:corpus});
\item \code{paper/figures/native\_nist.py} (Tables~\ref{tab:native} and~\ref{tab:vocab});
\item \code{paper/figures/nudge\_roundtrip.py} (Table~\ref{tab:nudge});
\item \code{paper/ir/tables.py}, which writes every table and inline number.
\end{itemize}
The software is archived as
\href{https://doi.org/10.5281/zenodo.22848721}{10.5281/zenodo.22848721}; the
companion recovery paper is
\href{https://doi.org/10.5281/zenodo.22848792}{10.5281/zenodo.22848792}.

\paragraph{How this was built.} The work was iterative: we steered the design
with Claude (Anthropic) over many rounds, pushing on correctness, generality and
performance; Claude wrote and debugged the code and ran the measurements.

\bibliographystyle{plainnat}
\bibliography{references}

\end{document}
